\chapter{Descrizione dello stage}
\label{cap:descrizione-stage}

\intro{Il presente capitolo descrive gli obiettivi del progetto di stage, la pianificazione delle attività in termini di ore di lavoro e l’analisi preventiva dei principali rischi che potrebbero insorgere nel corso del suo svolgimento.}\\

\section{Introduzione al progetto}

Il progetto di stage si è sviluppato lungo due linee di lavoro parallele, entrambe
rivolte al \emph{framework} \emph{Sia.Core}, base tecnologica condivisa da tutti i
prodotti aziendali. La prima linea ha riguardato la progettazione e l'implementazione
di un sistema di telemetria completo, basato sullo standard \gls{otel} e sullo stack
\emph{Grafana LGTM}, in grado di raccogliere e correlare \emph{trace}, metriche e
\emph{log} per offrire una visione unificata del comportamento delle applicazioni in
produzione. La seconda linea ha riguardato l'introduzione di un'infrastruttura di
\emph{testing} automatico nel \emph{framework}, articolata su \emph{unit test} e
\emph{integration test} lato \emph{server} e su \emph{unit test} e test
\emph{end-to-end} lato \emph{client}. Trasversalmente alle due linee, è stato adottato
un agente di intelligenza artificiale come \emph{pair engineer} strutturato, con la
creazione di agenti specializzati, di \emph{skill} procedurali e l'integrazione del
\gls{mcp} per consentire all'agente di interrogare direttamente i sistemi coinvolti
(\emph{Grafana}, database, strumenti di test). I fondamenti teorici dell'osservabilità
e della metodologia AI-assistita sono trattati in dettaglio rispettivamente nei
capitoli~\ref{cap:progettazione} e~\ref{cap:metodologia-ai}.

\section{Requisiti e obiettivi}

\subsection{Obiettivi}
\begin{enumerate}[label=\arabic*.]
  \item Analizzare i requisiti di osservabilità dell'ecosistema \emph{Sia.Core} e
        identificare le soluzioni tecnologiche adeguate, scegliendo lo stack di
        riferimento (\emph{OpenTelemetry} + \emph{Grafana LGTM}).
  \item Progettare un'architettura di telemetria basata sullo standard \gls{otel},
        integrandola con il \emph{framework} \emph{.NET 9} esistente in modo che ogni
        nuovo modulo possa aderirvi senza conoscerne i dettagli interni.
  \item Definire una tassonomia coerente di metriche, \emph{trace} e attributi per i
        moduli centrali del \emph{framework} (autenticazione, componenti \emph{Angular}
        come griglie e \emph{form}, prestazioni hardware, accesso al database),
        evitando l'esplosione della cardinalità lato Prometheus.
  \item Definire le strategie di campionamento delle \emph{trace} e di gestione delle
        metriche in ottica di scalabilità verso ambienti di produzione.
  \item Implementare meccanismi di arricchimento automatico delle \emph{trace} con
        l'identità dell'utente autenticato e di collegamento tra metriche aggregate e
        \emph{trace} tramite \emph{exemplar}.
  \item Configurare lo stack \emph{Grafana LGTM} e produrre dashboard in grado di
        rendere navigabile la catena diagnostica dal grafico alla singola richiesta.
  \item Introdurre da zero un'infrastruttura di \emph{testing} automatico nel
        \emph{framework} \emph{backend}, articolata su \emph{unit test} per la logica
        di business e \emph{integration test} sull'attendibilità dei dati nei
        database, sulle procedure SQL e sulle prospettive di configurazione dei
        componenti.
  \item Sviluppare test lato \emph{frontend} (\emph{Angular}), comprendenti
        \emph{unit test} per la logica applicativa e test di validazione
        \emph{end-to-end} per il comportamento dell'applicazione dal punto di vista
        dell'utente finale.
  \item Realizzare agenti AI specializzati, \emph{skill} procedurali e integrazioni
        \gls{mcp} per garantire la replicabilità della propagazione di telemetria e
        test ai diversi moduli del \emph{framework}.
\end{enumerate}

\subsection{Requisiti funzionali}

\begin{enumerate}[label=RF\arabic*.]
    \item Il sistema deve tracciare il percorso completo di ogni richiesta HTTP in
        ingresso, incluse le operazioni di accesso al database.
    \item Le operazioni CRUD verso il database devono essere riconoscibili tramite
        attributi che espongono il metodo e la \emph{repository} da cui è stata
        scatenata la richiesta.
    \item Il sistema deve esporre metriche operative per i moduli centrali del
        \emph{framework} (autenticazione, componenti \emph{Angular} come griglie e
        \emph{form}, prestazioni hardware, accesso al database).
    \item I \emph{log} devono essere correlabili alle \emph{trace} tramite
        \texttt{TraceId} e \texttt{SpanId}.
    \item Il sistema deve essere attivabile e disattivabile tramite configurazione,
        senza richiedere modifiche al codice.
    \item Deve essere possibile configurare indipendentemente il campionamento per
        ciascuna sorgente di \emph{trace}.
    \item Deve essere possibile visualizzare \emph{trace} e metriche raggruppate in
        dashboard \emph{Grafana} costruite sullo stack \emph{LGTM}.
    \item Le dashboard di \emph{Grafana} devono raccogliere i dati da più servizi e
        installazioni di clienti diversi, e permettere di filtrare per gli stessi.
    \item Il sistema deve arricchire automaticamente ogni \emph{span} con l'identità
        dell'utente autenticato (\texttt{enduser.id}, \texttt{enduser.name}) senza
        richiedere modifiche ai singoli moduli applicativi.
    \item Le metriche di tipo \emph{Histogram} devono esporre \emph{exemplar}, ovvero
        puntatori al \texttt{TraceId} del campione, per abilitare la navigazione dalla
        metrica aggregata alla \emph{trace} corrispondente.
    \item Devono essere prodotti \emph{unit test} per la logica di business dei moduli
        core del \emph{framework backend}, eseguibili in modo automatizzato.
    \item Devono essere prodotti \emph{integration test} per la verifica
        dell'attendibilità dei dati presenti nel database, delle procedure SQL e delle
        prospettive di configurazione dei componenti.
    \item Devono essere prodotti \emph{unit test} lato \emph{frontend} per la verifica
        della logica applicativa dei componenti \emph{Angular}.
    \item Devono essere prodotti test \emph{end-to-end} lato \emph{frontend} per
        verificare il corretto comportamento e la corretta visualizzazione
        dell'applicazione dal punto di vista dell'utente finale.
    \item La suite di test deve essere integrabile nel processo di \emph{build} e
        fornire un report di copertura del codice.
    \item Devono essere realizzati agenti AI specializzati e \emph{skill} procedurali
        per propagare in modo uniforme telemetria e test ai diversi moduli del
        \emph{framework}.
\end{enumerate}

\subsection{Requisiti non funzionali}

\begin{enumerate}[label=RNF\arabic*.]
    \item La telemetria non deve degradare sensibilmente le prestazioni in produzione.
    \item Le metriche devono essere progettate per non generare un numero eccessivo di serie temporali in Prometheus; in particolare, i tag ad alta variabilità (es.\ identificativi utente o di record) non devono essere usati come attributi di metriche.
    \item I segnali prodotti devono essere compatibili con qualsiasi backend \gls{otel}-compatibile.
    \item Le impostazioni devono essere modificabili per ambiente (sviluppo, \emph{staging}, produzione) tramite \emph{appsettings}.
\end{enumerate}


\section{Analisi preventiva dei rischi}
Durante la fase iniziale di analisi, sono stati individuati i principali rischi che sarebbero potuti emergere nel corso dello stage. 
Per ciascun rischio identificato, è stata definita una possibile strategia di mitigazione, al fine di ridurne l’impatto e garantire il corretto avanzamento delle attività previste.
\\
\begin{risk}{Tecnologie sconosciute}
    \riskdescription{alcune delle tecnologie previste nel progetto risultano totalmente o parzialmente sconosciute, con il rischio di rallentare le attività nelle fasi iniziali dello stage}
    \risksolution{pianificazione di un’attività strutturata di studio e approfondimento, attraverso l’utilizzo di documentazione ufficiale, tutorial online, corsi forniti dall’azienda e lo sviluppo di piccoli progetti di consolidamento, al fine di acquisire rapidamente le competenze necessarie}
    \label{risk:unknown-tecnologies} 
\end{risk}
\vspace{0.2cm}
\begin{risk}{Gestione del tempo}
    \riskdescription{la pianificazione delle attività potrebbe risultare non accurata, soprattutto a causa della presenza di tecnologie nuove e della componente sperimentale legata all’uso dell’intelligenza artificiale, con il rischio di ritardi rispetto alle scadenze previste}
    \risksolution{adottare il modello agile \gls{scrumg} per definire obiettivi intermedi settimanali, monitoraggio continuo dello stato di avanzamento e revisione periodica della pianificazione, in modo da adattare le attività in base alle difficoltà riscontrate}
    \label{risk:time-administration} 
\end{risk}
\vspace{0.2cm}
\begin{risk}{Allucinazioni \gls{llm}}
    \riskdescription{l'agente di codifica \emph{Claude Code}, utilizzato come \emph{pair engineer} per accelerare lo studio delle tecnologie e la propagazione della telemetria ai moduli del framework, può produrre risultati non corretti o non aderenti alle specifiche (fenomeno noto come ``allucinazione''), compromettendo l'affidabilità del codice generato}
    \risksolution{adozione di tecniche di verifica e validazione dei risultati prodotti, affiancate da un'attenta progettazione delle istruzioni fornite all'agente e dall'utilizzo di specifiche strutturate. Iterazioni successive e revisioni manuali del codice consentono di garantire la correttezza dell'instrumentazione introdotta nei moduli applicativi}
    \label{risk:LLM-hallucinations}
\end{risk}
\vspace{0.2cm}
\begin{risk}{Integrazione tra tecnologie}
    \riskdescription{l'integrazione tra \emph{OpenTelemetry}, il framework \emph{.NET 9} di \emph{Sia.Core} e la piattaforma di visualizzazione \emph{Grafana} (con i relativi backend Prometheus, Tempo e Loki) può presentare difficoltà legate a compatibilità, interfacciabilità e correlazione dei segnali, con il rischio di rallentare lo sviluppo o compromettere il corretto flusso dei dati telemetrici}
    \risksolution{utilizzo di un approccio incrementale all'integrazione, partendo da configurazioni di base del pipeline \gls{otel} e aumentando progressivamente la complessità. Consultazione della documentazione, utilizzo di ambienti di sviluppo locali per la validazione end-to-end dei segnali, e confronto con il team aziendale per la risoluzione di eventuali criticità}
    \label{risk:tecnologies-integration}
\end{risk}
\vspace{0.2cm}
\begin{risk}{Impatto della telemetria sulle prestazioni}
    \riskdescription{l'instrumentazione capillare dei moduli applicativi e la raccolta di trace, metriche e log potrebbero introdurre un overhead non trascurabile in produzione, degradando i tempi di risposta delle \gls{api} o generando un volume eccessivo di dati telemetrici}
    \risksolution{adozione di strategie di campionamento configurabili per sorgente, attenta progettazione degli attributi di metriche per evitare cardinalità elevata in Prometheus, e possibilità di disattivare la telemetria tramite \emph{appsettings} senza modifiche al codice. Validazione dell'overhead in ambiente di staging prima del rilascio in produzione}
    \label{risk:telemetry-overhead}
\end{risk}
\vspace{0.2cm}
\begin{risk}{Differenze di comportamento tra ambiente di sviluppo e produzione}
    \riskdescription{il sistema di telemetria potrebbe comportarsi in modo differente tra ambiente di sviluppo, \emph{staging} e produzione, a causa di differenze nella configurazione del collector \gls{otel}, nei volumi di traffico o nell'infrastruttura. Ciò può portare a malfunzionamenti non rilevati durante le fasi di sviluppo}
    \risksolution{utilizzo di ambienti di \emph{staging} il più possibile allineati alla produzione, configurazione differenziata per ambiente tramite \emph{appsettings}, e validazione end-to-end del flusso dall'evento applicativo alla dashboard \emph{Grafana} prima del rilascio}
    \label{risk:production-problems}
\end{risk}
\vspace{0.2cm}
\begin{risk}{Qualità dei test generati automaticamente}
    \riskdescription{i test generati tramite l'agente di codifica \emph{Claude Code} potrebbero risultare incompleti, non corretti o con copertura insufficiente, compromettendo l'efficacia del processo di \emph{testing}}
    \risksolution{introduzione di attività di validazione e revisione dei test generati, utilizzo di metriche di copertura del codice e iterazione sulle istruzioni fornite all'agente per migliorare progressivamente la qualità dei risultati}
    \label{risk:test-quality}
\end{risk}
\vspace{0.2cm}
\begin{risk}{Cardinalità eccessiva delle metriche}
    \riskdescription{l'utilizzo improprio di tag ad alta variabilità (es.\ identificativi utente, identificativi di record, parametri di query) come attributi di metriche potrebbe generare un numero esplosivo di serie temporali in Prometheus, compromettendo le prestazioni del backend di storage e l'usabilità delle dashboard}
    \risksolution{definizione di linee guida chiare sulla scelta degli attributi di metrica, separazione netta tra informazioni adatte a metriche aggregate e informazioni adatte a span e log, revisione delle istruzioni fornite all'agente \emph{Claude Code} per la propagazione della telemetria, e monitoraggio periodico della cardinalità in Prometheus}
    \label{risk:metrics-cardinality}
\end{risk}

\section{Pianificazione}
L’attività di stage è stata pianificata su una durata complessiva di 8 settimane, per un totale di 300 ore.
\\
La pianificazione è stata strutturata in modo progressivo, suddividendo il lavoro in fasi settimanali, ciascuna con obiettivi specifici e attività mirate, al fine di garantire un apprendimento graduale e un’integrazione efficace delle tecnologie coinvolte.
\subsection{Pianificazione settimanale}
\subsubsection{Settimane 1 e 2: Studio dell’architettura del sistema}
Le prime due settimane sono state dedicate all’analisi dell’architettura dell’applicazione Sia.Core, con particolare attenzione alla separazione tra componente \emph{client-side} e \emph{server-side}.
\begin{itemize}
    \item Analisi della struttura del progetto \emph{Angular}: moduli, componenti, servizi, \emph{routing}
    \item Studio delle convenzioni e pattern architetturali adottati (\emph{lazy loading, state management})
    \item Revisione del codice esistente con supporto del team di sviluppo
    \item Analisi della struttura del progetto \emph{.NET 9}: \emph{layer, \gls{api}, middleware, dependency injection}
    \item Studio dei pattern architetturali utilizzati (\emph{Clean Architecture, CQRS, repository pattern})
    \item Comprensione dei flussi di dati tra \emph{client} e \emph{server}: \emph{\gls{api} REST}, contratti, autenticazione
    \item Mappatura dei componenti critici candidati all'integrazione di test e telemetria
\end{itemize}

\subsubsection{Settimana 3: Studio degli strumenti e della metodologia AI-assistita}
La terza settimana è stata dedicata all'apprendimento degli strumenti chiave del progetto: lo stack \emph{OpenTelemetry} + \emph{Grafana LGTM}, gli strumenti di \emph{testing} per \emph{.NET} e \emph{Angular}, e la metodologia di lavoro basata su agenti di intelligenza artificiale.
\begin{itemize}
    \item \emph{OpenTelemetry}: studio dei concetti fondamentali (\emph{trace}, metriche, \emph{log}), del modello a \emph{span}, delle convenzioni semantiche e dell'integrazione con \emph{.NET 9}
    \item \emph{Grafana LGTM}: configurazione istanza locale, connessione ai backend \emph{Loki}, \emph{Grafana}, \emph{Tempo} e \emph{Mimir/Prometheus}, realizzazione di dashboard di base
    \item Studio degli strumenti di \emph{testing}: framework di \emph{unit test} e \emph{integration test} per \emph{.NET 9}, framework di \emph{unit test} e di test \emph{end-to-end} per \emph{Angular}
    \item Studio della metodologia AI-assistita: agenti specializzati, \emph{skill} procedurali e \gls{mcp} per l'interrogazione diretta di \emph{Grafana}, database e strumenti di test
    \item Realizzazione di un ambiente di sviluppo locale integrato con il \emph{pipeline} \gls{otel} completo
\end{itemize}

\subsubsection{Settimana 4: Progettazione dell'architettura di telemetria e strategia di testing}
Durante la quarta settimana è stata progettata l'architettura del sistema di telemetria ed è stata definita in parallelo la strategia di \emph{testing} per le due linee di lavoro.
\begin{itemize}
    \item Analisi dei moduli centrali del \emph{framework} candidati all'instrumentazione (autenticazione, componenti \emph{Angular} come griglie e \emph{form}, prestazioni hardware, accesso al database)
    \item Definizione del \emph{pipeline} \gls{otel}: \emph{ActivitySource}, \emph{Meter}, \emph{exporter} OTLP verso il \emph{collector}
    \item Progettazione del \emph{sampler} a due livelli e delle politiche di campionamento configurabili per sorgente
    \item Definizione di una tassonomia coerente di metriche, \emph{trace} e attributi, con regole esplicite per evitare l'esplosione della cardinalità lato Prometheus
    \item Identificazione delle tipologie di test: \emph{unit test} backend, \emph{integration test} su database, procedure SQL e prospettive, \emph{unit test} frontend, test \emph{end-to-end} frontend
    \item Definizione del flusso di lavoro AI-assistito (agenti specializzati, prompt strutturati, criteri di accettazione, revisione manuale)
    \item Stesura del documento di progettazione e della strategia di testing condivisi con il team
\end{itemize}

\subsubsection{Settimana 5: Implementazione del pipeline OTel e prima suite di test}
La quinta settimana è stata dedicata all'implementazione del \emph{pipeline} \emph{OpenTelemetry} nel \emph{framework} \emph{.NET 9}, alla prima instrumentazione dei moduli core e alla realizzazione di una prima suite di test con il supporto degli agenti AI.
\begin{itemize}
    \item Configurazione del \emph{pipeline} \gls{otel} nel \texttt{BaseApplicationBuilder} con \emph{sampler}, \emph{exporter} OTLP e instrumentazioni automatiche
    \item Integrazione di \emph{Serilog} con esportazione OTLP per la correlazione \emph{log}--\emph{trace} tramite \texttt{TraceId} e \texttt{SpanId}
    \item Implementazione del pattern \texttt{SiaXxxDiagnostics} per i moduli di autenticazione e accesso ai dati
    \item Arricchimento automatico degli \emph{span} con l'identità dell'utente autenticato (\texttt{enduser.id}, \texttt{enduser.name})
    \item Prima generazione di \emph{unit test} backend e \emph{unit test} frontend con il supporto degli agenti AI
    \item Validazione \emph{end-to-end} del flusso telemetrico e revisione manuale del codice di test generato
\end{itemize}

\subsubsection{Settimana 6: Propagazione della telemetria, dashboard e integration test}
Nella sesta settimana il lavoro si è concentrato sulla propagazione della telemetria a ulteriori moduli del \emph{framework}, sulla configurazione delle dashboard e sull'introduzione degli \emph{integration test} e dei test \emph{end-to-end} frontend.
\begin{itemize}
    \item Sviluppo di \emph{skill} procedurali per standardizzare la propagazione della telemetria ai diversi moduli del \emph{framework}
    \item Instrumentazione dei componenti \emph{Angular} (griglie, \emph{form}), del modulo prestazioni hardware e di ulteriori \gls{api} con \emph{trace} e metriche
    \item Definizione di metriche di tipo \emph{Histogram} con \emph{exemplar} per abilitare la navigazione dalla metrica aggregata alla \emph{trace} corrispondente
    \item Configurazione delle dashboard \emph{Grafana} con variabili \emph{template} per filtrare per servizio, installazione e cliente, sfruttando il \gls{mcp} per l'interrogazione diretta dello stack \emph{LGTM}
    \item Sviluppo di \emph{integration test} sull'attendibilità dei dati nel database, sulle procedure SQL e sulle prospettive di configurazione dei componenti
    \item Sviluppo dei test \emph{end-to-end} frontend per la verifica del comportamento dell'applicazione dal punto di vista dell'utente finale
\end{itemize}

\subsubsection{Settimana 7: Validazione del sistema integrato, tuning delle dashboard e suite di test}
La settima settimana ha previsto la validazione complessiva del sistema di telemetria su ambiente di \emph{staging}, l'ottimizzazione delle dashboard e l'esecuzione della suite completa di test.
\begin{itemize}
    \item Verifica della telemetria \emph{end-to-end}: dall'evento applicativo al \emph{collector} \gls{otel}, fino alla dashboard \emph{Grafana}
    \item Validazione dell'\emph{overhead} introdotto dalla telemetria sulle prestazioni applicative
    \item Verifica della cardinalità delle metriche in Prometheus e correzione di eventuali attributi problematici
    \item Ottimizzazione delle dashboard \emph{Grafana}: \emph{alert}, soglie, visualizzazioni avanzate e correlazione \emph{cross}-segnale (metrica $\rightarrow$ \emph{trace} via \emph{exemplar} $\rightarrow$ \emph{log} via \texttt{TraceId})
    \item Esecuzione della suite completa di test (\emph{unit}, \emph{integration} e \emph{E2E}) su ambiente di \emph{staging}, analisi della copertura del codice e correzione dei test falliti
    \item Stesura del report intermedio sullo stato del sistema di osservabilità e sulla qualità della suite di test
\end{itemize}

\subsubsection{Settimana 8: Messa in produzione e documentazione finale}
L'ultima settimana è stata dedicata alla messa in produzione della soluzione di telemetria, al consolidamento della suite di test e alla redazione della documentazione finale.
\begin{itemize}
    \item Preparazione e revisione finale dell'ambiente di produzione
    \item \emph{Deploy} della configurazione \emph{OpenTelemetry}, del \emph{collector} e delle dashboard \emph{Grafana LGTM} in produzione
    \item Consolidamento della suite di test e integrazione nel processo di \emph{build} con report di copertura del codice
    \item Stesura della documentazione tecnica finale (architettura del sistema di osservabilità, convenzioni di instrumentazione, dashboard, suite di test, agenti AI e \emph{skill} prodotti)
    \item Stesura della relazione di stage con analisi critica del lavoro svolto
    \item Presentazione finale dei risultati al team aziendale
\end{itemize}
